\documentclass[fontsize=11pt]{article}

\usepackage[margin=1in]{geometry}
\usepackage{graphicx}
\usepackage{wrapfig}
\usepackage{subcaption}
\usepackage{amsmath} 
\usepackage{siunitx}
\usepackage{booktabs}
\usepackage{gensymb}
\usepackage{hyperref}
\usepackage[export]{adjustbox}
\usepackage{anyfontsize}

\begin{document}

	\graphicspath{{./figures/}}

	%BJC: check address, phone number and fax (?) number
	\begin{figure}
	\centering
	\begin{minipage}{0.37\textwidth}
	\includegraphics[width=2.14in,left]{CUBoulder.pdf}
	\end{minipage}
	\begin{minipage}{0.62\textwidth}
	\scriptsize
	\noindent Department of Chemical and Biological Engineering \hfill t 303-735-7860~~~~~~~~~~~~~~~~~~ \\
	\noindent 596 UCB \hfill f 303-492-8425~~~~~~~~~~~~~~~~~~ \\
	\noindent Boulder, Colorado 80309 \hfill michael.shirts@colorado.edu \\
	\end{minipage}
	\end{figure}
	
	%\noindent May 3\textsuperscript{rd}, 2019\\
	\noindent \today \\

	\noindent Dear Editors,\\
	
	\newcommand{\ManuscriptTitle}{Capturing Subdiffusive Solute Dynamics and 
	Predicting Selectivity in Nanoscale Pores with Time Series Modeling}
	
	We are submitting ``\ManuscriptTitle'' by Benjamin J. Coscia and Michael 
	R.  Shirts, for consideration for publication as an Article in the Journal
	of Chemical Theory and Computation. We expect this work to be of significant
	interest to researchers who study complex transport at the molecular level,
	as it describes important new approaches for estimating long time scale 
	behavior and bulk transport behavior from shorter molecular trajectories 
	exhibiting anomalous diffusion rather than Brownian motion. Please note 
	that this article has been uploaded to arXiv with differences only in 
	formatting: \href{https://arxiv.org/abs/2004.07905}{https://arxiv.org/abs/2004.07905}. \\        
%        Please note that the content of this article
%        %BJC: see if you agree with these changes
%        has been uploaded to arXiv: \href{https://arxiv.org/abs/2004.07905}{https://arxiv.org/abs/2004.07905}.
%        For submission to JCTC, we have made minor changes to the introduction,
%        abstract and conclusion in order to make it more suited for readers of 
%        the journal. We also made efforts to shorten the manuscript by moving 
%        certain discussions and figures to the supporting information. \\
	
    In this work, we build on the chemical intuition gained from our previous 
    study (Coscia and Shirts \textit{J. Phys.  Chem. B}, 123, 6314--6330 (2019))
    in order to formulate two different stochastic time series modeling approaches,
    %which produce trajectory realizations that share the same dynamical properties as 
    %those exhibited by solutes in our 5~$\mu$s molecular dynamics (MD) simulations.
    parameterized from 5~$\mu$s molecular dynamics (MD) simulations, capable of 
    predicting solute flux and selectivity on a macroscopic scale. 
    First, using anomalous diffusion theory, we parameterize solute motion as 
    subordinated fractional Brownian and L\'evy motion processes, based on the 
    observation that solutes experience a series of dwell times separated by
    anti-correlated hop lengths. Secondly, we add state-dependent dynamics to 
    standard Markov state modeling approaches, where we model solute motion as 
    a series of transitions between various trapped states, each with different
    dynamical properties. We demonstrate how one can project these models 
    onto macroscopic timescales in order to estimate solute flux and selectivity.
    This is among the first studies to apply time series modeling of molecular 
    simulations in order to infer behavior at long time scales and estimate bulk
    materials properties with complex anomalous diffusion mechanisms.
	
%	In this article, we analyze solute trajectories generated from 5~$\mu$s molecular 
%	dynamics (MD) simulations of an inverted hexagonal (H\textsubscript{II}) phase self-assembled lyotropic
%        liquid crystal (LLC) polymer membrane to parameterize two
%        different stochastic time series approaches capable of
%        predicting solute flux and selectivity on a macroscopic
%        scale. Both approaches quantitatively reproduce the mean
%        squared displacements (MSDs) exhibited by solutes in MD
%        simulations. This is among the first studies to apply time
%        series modeling of molecular simulations in order to infer
%        behavior at long time scales and estimate bulk
%        materials properties with complex anomalous diffusion
%        mechanisms.\\
%
%	This article builds on recent work of ours published in
%        J. Phys. Chem. B where we characterized transport mechanisms
%        exhibited by a series of small polar solutes in an LLC
%        Membrane (Coscia and Shirts \textit{J. Phys.  Chem. B}, 123,
%        6314--6330 (2019)).  We explore two main modeling
%        approaches. First, using anomalous diffusion theory, we
%        parameterize solute motion as subordinated fractional Brownian
%        and L\'evy motion processes, based on the observation that
%        solutes experience a series of dwell times separated by
%        anti-correlated hop lengths. Secondly, we add state-dependent
%        dynamics to standard Markov state modeling approaches, where
%        we model solute motion as a series of transitions between
%        various trapped states, each with different dynamical
%        properties.  Finally, we demonstrate how one can project models 
%        onto macroscopic timescales in order to estimate
%        solute flux and selectivity.\\
	
	%BJC: included phone numbers because it asked to do so in author's guide.
	\noindent Some suggestions for reviewers are:
	\begin{enumerate}
	
		\item Francisco Hung has used molecular simulations to study nematic liquid crystals
		and transport under confinement within various nanostructured materials 
		(Northeastern University, 617-373-8619,\\ \href{mailto:f.hung@northeastern.edu}{f.hung@northeastern.edu}).
		
		\item Mahesh Mahanthappa has experience modeling soft materials using molecular dynamics
		including work on ion transport in various ordered liquid crystalline phases. 
		(University of Minnesota, 612-625-4599, \href{mailto:maheshkm@umn.edu}{maheshkm@umn.edu}).
		
		\item Stephen Paddison uses various levels of molecular simulation theory to understand
		proton transport in nanostructured and charged polymer membranes. (The University of Tennessee 
		Knoxville, 865-974-2026, \href{mailto:spaddison@utk.edu}{spaddison@utk.edu}).
		
		\item Coray Colina uses various simulation techniques to understand structure-property relations
		in functional materials including nanoporous materials for separation applications. (University
		of Florida, 352-294-3488, \href{mailto:colina@chem.ufl.edu}{colina@chem.ufl.edu}).
		
		\item Venkat Ganesan uses multiscale modeling to study solute and ion transport in complex
		nanostructured morphologies in order to learn how to design novel materials. (The University
		of Texas at Austin, 512-471-4856, \href{mailto:venkat@che.utexas.edu}{venkat@che.utexas.edu}).
		
%		\item Ralf Metzler is an expert in anomalous stochastic processes, with interest in the physical
%		origins of anomalous diffusion as well as the inference of the parameters which govern these
%		processes. (Universität Potsdam, +49 331 977 5715, \href{mailto:rmetzler@uni-potsdam.de}{rmetzler@uni-potsdam.de}).
		
	\end{enumerate}
	
	\noindent Please send correspondence regarding this paper to Michael R. Shirts (contact
	details in letterhead).\\	
	
	\noindent Sincerely,
	
	\noindent Michael R. Shirts \\
	\noindent Benjamin J. Coscia \\
	
\end{document}

% LocalWords:  BJC Subdiffusive Solute Nanoscale Coscia arXiv solute lyotropic
% LocalWords:  LLC MSDs solutes Phys Chem subdiffusive evy timescales Klafter
% LocalWords:  univeristy JPCB woould nanostructured Mahesh Mahanthappa klafter
% LocalWords:  Tel Paddison
